% Part: methods
% Chapter: proofs
% Section: starting-proofs

\documentclass[../../../include/open-logic-section]{subfiles}

\begin{document}

\olfileid{mth}{prf}{str}

\olsection{كيف يُبتدأ البرهان}

فمن أي موضع يكون الابتداء؟

المطلوب الذي أُعطيته لتثبته هو منتهى البرهان، فلا تجعله من مقدماته.
ولك أن \emph{تعلن} في أوله ما تريد إثباته، لكن ليس لك أن تتخذه
فرضًا. ومن النافع أن تكتب المطلوب في أسفل ورقة خالية، ليظل المقصد
نصب عينيك وأنت ترتب الطريق إليه.

ثم انظر فيما أُبيح لك من الفروض؛ وسيظهر أثر \emph{نوع} البرهان في
ذلك ظهورًا أوفى في القسم الآتي. أثبت الفروض في صدر الصفحة وصرّح
بكونها فروضًا: فإن كان~${\OLId{latin-ordinary}{p}}$ مفروضًا فقل «افترض أن~${\OLId{latin-ordinary}{p}}$» أو «لنفترض
أن~${\OLId{latin-ordinary}{p}}$». ثم استخرج من السؤال كل تعريف تحتاج إليه؛ فقد ينص السؤال على
تعريف بعينه، أو تشتمل فروضه أو نتيجته على ألفاظ معرّفة متعددة.
\emph{فاكتب تلك التعريفات، وتحقق من فهم مدلولها.}

ولصورة السؤال أيضًا حكم في إعداد البرهان؛ لذلك يفصل القسم التالي طرق
الإعداد بحسب نوع الجملة المراد إثباتها.

\end{document}
